\documentclass[../main.tex]{subfiles}

\begin{document}

\chapter{一般测度空间上的积分}\label{chapter:ch18-integration}

我们从 18.1 节考虑可测函数开始对一般测度空间上的积分进行研究, 许多内容与一元实变量的 Lebesgue 可测函数的研究十分类似. 我们研究一般积分的途径与第 4 章中探求的关于一元实变量函数的 Lebesgue 测度的积分不同. 在 18.2 节, 我们首先对非负简单函数定义积分, 接着直接定义非负可测函数 $f$ 的积分为满足 $0\leqslant\psi\leqslant f$ 的非负简单函数 $\psi$ 的积分的上确界. 在这个早期阶段我们证明一般的 Fatou 引理, 它是积分理论完全发展的基石, 接下来证明与它密切相关的结论——单调收敛定理与 Beppo Levi 引理. 在 18.3 节, 我们考虑一般可测函数的积分且证明该积分的线性与单调性、连续性、可数可加性, 以及积分比较判别法、Vitali 收敛定理. 在 18.4 节, 我们引入一个测度关于另一个测度绝对连续的概念且证明 Radon-Nikodym 定理, 它是将一元实变量的绝对连续函数表示为不定积分这一结果的深远推广. 我们也证明了关于测度的 Lebesgue 分解定理. 在本章最后, 我们用 Baire 范畴定理证明 Vitali-Hahn-Nikodym 定理, 它告诉我们在非常一般的假设下, 测度序列的集合态极限仍然是一个测度.

\subfile{Chapter-18/section-01.tex}

\subfile{Chapter-18/section-02.tex}

\subfile{Chapter-18/section-03.tex}

\subfile{Chapter-18/section-04.tex}

\subfile{Chapter-18/section-05.tex}

\end{document}